\documentclass[a4paper,11pt]{article}
\usepackage{fontspec}
\setmainfont{Noto Sans CJK JP}
\setsansfont{Noto Sans CJK JP}
\setmonofont{Noto Sans Mono CJK JP}

% CJK line breaking (XeTeX primitives)
\XeTeXlinebreaklocale "ja"
\XeTeXlinebreakskip = 0pt plus 1pt minus 0.1pt
\XeTeXlinebreakpenalty = 0

\usepackage{geometry}
\geometry{margin=2.5cm}

\usepackage{xcolor}
\definecolor{accent}{HTML}{2980B9}
\definecolor{darkblue}{HTML}{1A3A5F}
\definecolor{codebg}{HTML}{F5F5F5}
\definecolor{codeframe}{HTML}{DDDDDD}
\definecolor{infoframe}{HTML}{2980B9}
\definecolor{infobg}{HTML}{EBF5FB}

\usepackage{fancyhdr}
\pagestyle{fancy}
\fancyhead[L]{}
\fancyhead[R]{\small\color{accent}runops + Claude Code 導入ガイド}
\fancyfoot[C]{\thepage}
\renewcommand{\headrulewidth}{0.4pt}

\usepackage{listings}
\lstset{
  basicstyle=\small\ttfamily,
  backgroundcolor=\color{codebg},
  frame=single,
  rulecolor=\color{codeframe},
  breaklines=true,
  columns=fullflexible,
  keepspaces=true,
  xleftmargin=4pt,
  xrightmargin=4pt,
  framesep=4pt,
  framexleftmargin=0pt,
  commentstyle=\color{green!40!black}\itshape,
  language=bash,
}

\usepackage{fancybox}
\usepackage{booktabs}
\usepackage{longtable}
\usepackage{enumitem}
\usepackage{hyperref}
\hypersetup{colorlinks=true, linkcolor=darkblue, urlcolor=accent}

\renewcommand{\arraystretch}{1.3}
\setlength{\parindent}{0pt}
\setlength{\parskip}{6pt plus 2pt minus 1pt}

% Info box
\newenvironment{infobox}{%
  \medskip\noindent
  \fboxrule=0.5pt\fboxsep=8pt
  \fcolorbox{infoframe}{infobg}\bgroup\minipage{\dimexpr\linewidth-2\fboxsep-2\fboxrule}
}{%
  \endminipage\egroup\medskip
}

\begin{document}

% ====== Title page ======
\thispagestyle{empty}
\vspace*{5cm}

\begin{center}
{\Huge\bfseries\color{darkblue} runops + Claude Code}

\vspace{1cm}

{\Large\color{accent} HPC シミュレーション研究のための\\[4pt] AI 支援ワークフロー}

\vspace{2cm}

{\color{accent}\rule{8cm}{0.5pt}}

\vspace{1.5cm}

{\large 対象環境: camphor3 (京都大学) + EMSES}

\medskip

runops v0.3.0 ~/~ Claude Code (Anthropic)

\vspace{4cm}

{\small\color{gray} 2026-04-10}
\end{center}

\newpage

% ====== Section 1 ======
\section*{\color{darkblue} 1.\quad はじめに --- なぜ runops + Claude Code か}
\addcontentsline{toc}{section}{はじめに}

HPC シミュレーション研究では、パラメータファイルの設定、run の準備・投入・管理、解析・可視化、考察のまとめといった定型作業に多くの時間を取られます。これらは研究に不可欠ですが、研究者が本来やりたいこと --- 物理の考察、新しい仮説の検証、論文の執筆 --- ではありません。

また、パラメータ設定のコピペミスは計算資源の浪費に直結します。手作業が多いほどミスが増え、数日分の計算が無駄になることもあります。

\textbf{runops} は、これらの定型作業を構造化し AI Agent に委譲するために設計された HPC シミュレーション管理ツールです。\textbf{Claude Code} (Anthropic 社の AI コーディングエージェント) と組み合わせることで、\textbf{「やるべきこと」ではなく「やりたいこと」に集中できる}ワークフローを実現します。

さらに、すべての run が構造化・記録されるため、過去の実験が将来の研究資産として再利用可能になります。

\begin{infobox}
\textbf{大前提: ユーザーが覚えるべきは「設計思想」だけ。}\\[4pt]
runops のコマンド体系やオプションを覚える必要はありません。
日常の操作 --- ケース作成、パラメータサーベイ設計、run 生成・投入、状態監視、解析、知見整理、Git コミット --- はすべて \textbf{Claude Code (AI Agent) に自然言語で指示} するだけで実行されます。ユーザーが直接行うのは \textbf{(1) 初期セットアップ} と \textbf{(2) ジョブ投入などの承認操作} だけです。本資料では runops の設計思想と AI Agent の活用法を中心に解説します。
\end{infobox}

% ====== Section 2 ======
\section*{\color{darkblue} 2.\quad runops の概要}

runops は Python 製の CLI ツールで、以下の機能を提供します。

\subsection*{\color{accent} 主な機能}

\begin{itemize}[nosep]
  \item \textbf{Run 中心の管理:} 各シミュレーション実行を一意な run\_id で識別し、入力・出力・設定・来歴を 1 つのディレクトリに集約
  \item \textbf{パラメータサーベイ:} survey.toml にサーベイ軸を定義し、直積展開で run を自動生成 (手作業でのパラメータ計算が不要)
  \item \textbf{Slurm 統合:} ジョブスクリプトの自動生成、投入 (sbatch)、状態同期 (squeue/sacct)、ログ取得を一元管理
  \item \textbf{Simulator Adapter:} EMSES・BEACH 等のシミュレータごとに入力形式・実行方法・出力検出を抽象化 (新シミュレータの追加が容易)
  \item \textbf{知見管理:} 実験の知見を curated knowledge (insights + facts) と lab notebook (時系列ノート) の 2 層で記録
  \item \textbf{来歴 (Provenance):} manifest.toml に run の全情報を記録し、「いつ・誰が・どの設定で・どの commit から」実行したか追跡可能
\end{itemize}

\subsection*{\color{accent} 基本概念: 4 層の階層構造}

runops はシミュレーション研究を 4 つの階層で整理します:

\begin{center}
\begin{tabular}{lll}
\toprule
\textbf{階層} & \textbf{定義ファイル} & \textbf{役割} \\
\midrule
Campaign & campaign.toml & 研究テーマ全体 (仮説・変数・観測量) \\
Case & case.toml & シミュレーション設定のテンプレート \\
Survey & survey.toml & パラメータサーベイ定義 (掃引軸) \\
Run & manifest.toml & 個々のシミュレーション実行 (自動生成) \\
\bottomrule
\end{tabular}
\end{center}

例えば「イオン温度を 1--19\,eV で 10 点掃引」という実験は、Campaign の中の 1 つの Survey として定義され、AI Agent が自動的に 10 個の Run を生成します。各 Run は固有の run\_id を持ち、入力ファイル・ジョブスクリプト・manifest が自動で配置されます。

\begin{infobox}
これらの階層構造を理解しておけば、実際のファイル操作やコマンド実行は AI Agent が行います。ユーザーは「Ti を 1--19\,eV で 10 点掃引したい」と自然言語で伝えるだけで、AI が campaign.toml から survey.toml、run 生成まで一貫して処理します。
\end{infobox}

% ====== Section 3 ======
\section*{\color{darkblue} 3.\quad 初期セットアップ (唯一の手動操作)}

ユーザーが直接コマンドを実行するのは、このセットアップだけです。以降の操作はすべて Claude Code 経由で行います。camphor のような HPC 環境では sudo が使えないため、すべてユーザー権限でインストールします。

\subsection*{\color{accent} 前提条件}

\begin{itemize}[nosep]
  \item camphor3 のアカウント (Slurm が使える環境)
  \item Python 3.10 以上 (\texttt{module load} で利用可能)
  \item Git / GitHub アカウント
\end{itemize}

\subsection*{\color{accent} Step 1: uv のインストール}

Python パッケージマネージャ。runops のインストールに必要です:

\begin{lstlisting}
curl -LsSf https://astral.sh/uv/install.sh | sh
source ~/.bashrc   # パスを反映
\end{lstlisting}

\subsection*{\color{accent} Step 2: Node.js のインストール (nvm 経由)}

Claude Code と gh (GitHub CLI) の実行に Node.js が必要です。sudo 不要な nvm を使います:

\begin{lstlisting}
curl -o- https://raw.githubusercontent.com/nvm-sh/nvm/v0.40.3/install.sh | bash
source ~/.bashrc   # nvm を反映
nvm install --lts  # 最新 LTS 版をインストール
\end{lstlisting}

\subsection*{\color{accent} Step 3: gh (GitHub CLI) のインストールとログイン}

AI Agent がフィードバック issue を投稿するのに gh を使います。Node.js 環境があれば npm でインストールできます:

\begin{lstlisting}
npm install -g @anthropic-ai/claude-code  # Claude Code (後述)
npm install -g gh                         # GitHub CLI (非公式npm)
\end{lstlisting}

npm に gh パッケージがない場合は、公式バイナリを直接ダウンロードします:

\begin{lstlisting}
# Linux x86_64 の場合
GH_VER=2.74.0  # https://github.com/cli/cli/releases で最新版を確認
curl -L "https://github.com/cli/cli/releases/download/\
v${GH_VER}/gh_${GH_VER}_linux_amd64.tar.gz" | \
  tar xz --strip-components=1 -C ~/.local
\end{lstlisting}

インストール後、GitHub にログインします:

\begin{lstlisting}
gh auth login
\end{lstlisting}

対話プロンプトに従い、\texttt{GitHub.com} → \texttt{HTTPS} → ブラウザ認証 (または token 入力) を選択します。camphor からブラウザが開けない場合は \texttt{Paste an authentication token} を選び、\texttt{https://github.com/settings/tokens} で生成した Personal Access Token を貼り付けてください。

\subsection*{\color{accent} Step 4: runops でプロジェクトを初期化}

\texttt{uvx} (uv の一時実行機能) を使い、1 コマンドで初期化できます:

\begin{lstlisting}
mkdir my_simulation && cd my_simulation && git init
uvx --from runops runo init
source .venv/bin/activate
\end{lstlisting}

これにより runops のインストール、venv 作成、プロジェクト構造 (runops.toml, simulators.toml, campaign.toml, .runops/ 等) の生成がすべて自動で行われます。

\subsection*{\color{accent} Step 5: Claude Code を起動 → 以降は AI に任せる}

\begin{lstlisting}
claude
\end{lstlisting}

起動後は自然言語で指示するだけです:

\begin{lstlisting}[language={},commentstyle={}]
> 環境を確認して (runo doctor を実行)
> EMSES 用の flat_plate ケースを作って
> Ti=1-19eV で 10 点のサーベイを設計して
> smoke test として 3 点を gr20001a に投入して
> 完了した run を解析して
\end{lstlisting}

ケース作成、サーベイ設計、run 生成・投入、監視、解析、知見記録 --- これらはすべて AI Agent が runops コマンドを適切に選択・実行します。

% ====== Section 4 ======
\section*{\color{darkblue} 4.\quad Claude Code との連携}

\subsection*{\color{accent} Claude Code とは}

Claude Code は Anthropic 社が提供する AI コーディングエージェントです。ターミナル上で対話的に動作し、ファイルの読み書き・コマンド実行・Git 操作などを自律的に行えます。

\begin{itemize}[nosep]
  \item ターミナルで \texttt{claude} コマンドを起動するだけで使える CLI ツール
  \item コードの読解・編集・デバッグに加え、ファイル操作やシェルコマンドの実行が可能
  \item プロジェクトの CLAUDE.md を読んで文脈を理解し、プロジェクト固有のルールに従う
  \item 危険な操作 (ジョブ投入、ファイル削除等) は人間の承認を求める安全設計
\end{itemize}

\subsection*{\color{accent} runops のハーネス (Agent 連携基盤)}

runops は Claude Code 用のハーネスを自動生成します。\texttt{runo update-harness} コマンドにより以下が配置されます:

\begin{center}
\begin{tabular}{ll}
\toprule
\textbf{ファイル} & \textbf{役割} \\
\midrule
CLAUDE.md & プロジェクト運用ルール・スキル一覧・役割分担 \\
.claude/skills/ & 定型作業のスキル定義 (15+ 種類) \\
.claude/rules/ & Agent が守るべき制約 (ファイル保護、承認フロー等) \\
.claude/settings.json & 権限設定 (許可/確認/拒否するツール) \\
\bottomrule
\end{tabular}
\end{center}

\subsection*{\color{accent} 自然言語で指示 → AI が自動で実行}

runops のハーネスには「スキル」と呼ばれる定型作業の手順書が内蔵されており、AI Agent はユーザーの自然言語の指示に応じて適切なスキルを自動選択・実行します。\textbf{ユーザーがスキル名を覚える必要はありません。}

\begin{center}
\begin{tabular}{ll}
\toprule
\textbf{自然言語の指示例} & \textbf{AI が行うこと} \\
\midrule
「EMSES 用の新しいケースを作って」 & ケース作成 \\
「Ti を 10 点掃引するサーベイを設計して」 & パラメータサーベイ設計 \\
「全 run を生成して投入して」 & run 生成 → 投入 (承認あり) \\
「今の進捗を確認して」 & 状態確認・Slurm 同期 \\
「完了した run を解析して」 & 解析・集計 \\
「失敗した run を調べて」 & 失敗 run の診断 \\
「結果を知見として記録して」 & 知見の整理・保存 \\
「今の作業をノートに残して」 & lab notebook に記録 \\
\bottomrule
\end{tabular}
\end{center}

\begin{infobox}
\textbf{具体例:}\\
「Ti=1--19eV で 10 点サーベイを設計して、smoke test 3 点を gr20001a に投入して」\\
→ AI がサーベイ設計 → run 生成 → 投入提案 (ユーザー承認) → 状態確認まで一貫して実行
\end{infobox}

\subsection*{\color{accent} 安全設計: 人間の承認が必要な操作}

以下の操作は AI が自動実行せず、必ずユーザーの確認を求めます:

\begin{itemize}[nosep]
  \item \textbf{runo runs submit} (ジョブ投入): 対象 run・queue・資源量を提示してから実行
  \item \textbf{runo runs delete / purge-work} (実ファイル削除)
  \item 大規模 survey の\textbf{一括投入} (\texttt{--all})
  \item 設定ファイル (CLAUDE.md, settings.json 等) の変更
\end{itemize}

% ====== Section 5 ======
\section*{\color{darkblue} 5.\quad 典型的なワークフロー例}

実際の研究サイクルを例に、runops + Claude Code のワークフローを示します。

\subsection*{\color{accent} Phase 1: 計画 (人間 → AI に伝える)}

研究の方向性を AI Agent に自然言語で伝えます。AI が campaign.toml の作成・編集を行います。

\begin{lstlisting}[language={},commentstyle={}]
> 研究テーマは「太陽風プラズマ中の吸収平板によるイオン枯渇」。
  仮説は「枯渇角はイオン温度に依存する」。
  Ti を 1-19eV で掃引して、イオン密度と電位を観測したい。
  ベース入力は ~/large1/exp_flat/plasma.inp を使って。
\end{lstlisting}

\subsection*{\color{accent} Phase 2: 準備 (AI が自動実行)}

AI がケース作成、サーベイ設計、run 生成を一貫して行い、設計の経緯を lab notebook に自動記録します。

\begin{lstlisting}[language={},commentstyle={}]
> Ti サーベイを設計して、run を生成して
\end{lstlisting}

\subsection*{\color{accent} Phase 3: 実行 (AI 提案 → 人間承認)}

ジョブ投入は AI が提案し、\textbf{ユーザーが承認}してから実行されます。投入後の状態監視も AI に任せられます。

\begin{lstlisting}[language={},commentstyle={}]
> smoke test 3 点を gr20001a に投入して
> [AI が投入内容を提示 -> ユーザーが承認]
> 進捗を確認して
\end{lstlisting}

\subsection*{\color{accent} Phase 4: 解析 (AI が自動実行)}

\begin{lstlisting}[language={},commentstyle={}]
> 完了した run を解析して、Ti 依存性をプロットして
> run 間で比較して
\end{lstlisting}

\subsection*{\color{accent} Phase 5: 知見整理 (AI が自動実行)}

\begin{lstlisting}[language={},commentstyle={}]
> 結果を知見として記録して
> 次のサーベイの方針を議論しよう
\end{lstlisting}

AI が curated knowledge (insights + facts) を保存し、次のサーベイの方針を提案します。ユーザーとの議論を経て Phase 1 に戻ります。

% ====== Section 6 ======
\section*{\color{darkblue} 6.\quad 運用フィードバックの仕組み}

runops は開発中のツールであり、日常運用の中で改善点が見つかることがあります。runops にはこのフィードバックループを組織化するための仕組みが組み込まれています。

\subsection*{\color{accent} upstream-feedback ルール}

runops のハーネスには、AI Agent が日常運用中に runops 自体の改善点 (バグ・使いにくさ・不足機能・わかりにくいエラーメッセージなど) を検出したとき、GitHub issue としてフィードバックを提出するルールが組み込まれています。AI が自動で気づくこともありますが、ユーザーから「この挙動は runops に issue を出したほうがいい」と指示することもできます:

\begin{enumerate}[nosep]
  \item \textbf{AI が改善提案をユーザーに説明} (勝手に issue を切らない)
  \item ユーザーが OK したら \textbf{AI が gh issue create で issue を作成}
  \item 投げた issue の URL を \textbf{lab notebook に自動記録}
  \item 当面の研究作業は \textbf{workaround で継続} (フィードバックで作業を止めない)
\end{enumerate}

\begin{infobox}
\textbf{実例:} 実際の ion\_depression プロジェクトでは、運用中に 5 件の改善提案が GitHub issue として提出され、うち数件は v0.3.0 で対応済みです。例: encoding 修正、\texttt{--qos} オプション追加、\texttt{runs sync} の集約表示など。
\end{infobox}

\subsection*{\color{accent} /update-runops スキル (ユーザーが手動で呼ぶ)}

runops 本体やハーネスファイルの更新は、ユーザーが Claude Code 上で \texttt{/update-runops} と入力して実行します。これは意図的なタイミングで行う操作のため、自動ではなくスキルとして明示的に呼び出します:

\begin{lstlisting}[language={},commentstyle={}]
> /update-runops
AI: tools/runops を pull して update-harness を実行します...
    -> CLAUDE.md, skills/, rules/ が最新テンプレートで更新されました
    -> runo doctor で環境を確認 ... OK
\end{lstlisting}

実行内容:
\begin{itemize}[nosep]
  \item \textbf{tools/runops/} リポジトリを \texttt{git pull} で最新化
  \item \textbf{runo update-harness} でハーネスファイル (CLAUDE.md, skills, rules 等) を再生成
  \item \textbf{runo update} でシミュレータパッケージを更新
  \item 編集済みファイルとの競合は \texttt{.new} ファイルとして出力され、diff 確認後にマージ
\end{itemize}

チームで runops を使う場合、開発者が runops 本体を更新したら、各プロジェクトで \texttt{/update-runops} を実行するだけでハーネスが最新化されます。

\subsection*{\color{accent} フィードバックが runops を育てる}

重要なのは、「この 1 回だけ回避できればよい」ではなく、「次に同じ作業をする人が困らないようにする」という視点です。AI が運用中に自然とフィードバックを収集し、ユーザーの承認を経て upstream に還元する --- このサイクルにより runops は\textbf{使いながら育つツール}になります。

% ====== Section 7 ======
\section*{\color{darkblue} 7.\quad Claude Code の補足}

\subsection*{\color{accent} 認証}

Claude Code の初回起動時に Anthropic の認証が求められます。Max プラン (月額サブスクリプション) の場合は API キーなしでブラウザ認証で直接利用可能です。API プランの場合は \texttt{https://console.anthropic.com/} で API キーを取得してください。

\subsection*{\color{accent} 使い方のコツ}

\begin{itemize}[nosep]
  \item \textbf{具体的に指示する} --- 「解析して」より「完了した run の Ti 依存性をプロットして」のほうが精度が高い
  \item \textbf{承認プロンプトを確認する} --- ジョブ投入やファイル削除など危険な操作では AI が確認を求めるので、内容を読んでから承認する
  \item \textbf{会話を続ける} --- 1 つのセッション内で計画→準備→解析まで一貫して指示できる
\end{itemize}

\begin{infobox}
Claude Code は CLAUDE.md を自動的に読み込むため、runops プロジェクトのルール・スキル・制約を最初から理解した状態で動作します。特別な設定は不要です。
\end{infobox}

% ====== Section 8 ======
\section*{\color{darkblue} 8.\quad 運用のコツとベストプラクティス}

\subsection*{\color{accent} Git コミットは AI が自動で行う}

runops のハーネスには「意味のある作業単位ごとに Git コミットする」というルールが組み込まれており、AI Agent は case 作成後・run 生成後・解析完了後・知見保存後などのタイミングで自動的にコミットを作成します。もしコミットされていないと感じたら、「コミットして」と指示すれば即座に対応します。

\subsection*{\color{accent} Lab notebook は自然言語で指示する}

AI Agent に「今の議論をノートに残して」「サーベイの設計理由を記録して」と伝えれば、lab notebook (\texttt{notes/YYYY-MM-DD.md}) に自動記録します。「なぜこのパラメータ範囲にしたか」「なぜこの case を保留したか」といった情報は再現性に直結するので、積極的にメモを指示してください。\texttt{/note} スキルを直接呼ぶこともできますが、自然言語での指示のほうが柔軟です。

\subsection*{\color{accent} Smoke test を活用する}

大規模サーベイの前に、両端と中央の 2--3 点だけ先行投入 (smoke test) して入力の妥当性を確認するのが効果的です。「まず 3 点だけ smoke test して」と AI に指示すれば、適切な点を選んで投入します。

\subsection*{\color{accent} プロジェクトを GitHub private repo で管理する}

プロジェクトディレクトリは GitHub の private リポジトリとして管理することを強く推奨します。runops はシミュレーションの出力データ (大容量) を \texttt{.gitignore} で除外し、パラメータファイル・設定・解析スクリプト・知見など再現に必要な情報だけを Git 管理するように設計されているため、GitHub での管理が容易です。

これにより:
\begin{itemize}[nosep]
  \item HPC 上のデータ消失に対するバックアップ
  \item 研究の全履歴 (設計意図、パラメータ変更、解析結果) の永続的な記録
  \item 将来の自分や共同研究者が過去の実験を再現・参照可能
\end{itemize}

リポジトリの作成も AI Agent に任せられます:

\begin{lstlisting}[language={},commentstyle={}]
> このプロジェクトを GitHub に private repo として上げて
\end{lstlisting}

AI が \texttt{gh repo create} を実行し、リモート設定と初回 push まで行います。

\subsection*{\color{accent} バグを見つけたら feedback を}

runops のバグや改善点を見つけたら、AI Agent に「この挙動は runops に issue を出して」と伝えてください。AI がユーザーの承認を得てから GitHub issue を作成し、当面の研究は workaround で継続します。

% ====== Section 9 ======
\section*{\color{darkblue} 9.\quad 導入チェックリスト}

runops + Claude Code を始めるためのステップ:

\begin{enumerate}[nosep]
  \item \textbf{uv} --- \texttt{curl -LsSf https://astral.sh/uv/install.sh | sh}
  \item \textbf{nvm + Node.js} --- \texttt{curl -o- https://\ldots/nvm/install.sh | bash \&\& nvm install --lts}
  \item \textbf{Claude Code + gh} --- \texttt{npm install -g @anthropic-ai/claude-code gh}
  \item \textbf{gh ログイン} --- \texttt{gh auth login}
  \item \textbf{プロジェクト初期化} --- \texttt{mkdir my\_sim \&\& cd my\_sim \&\& git init}
  \item \textbf{runops 初期化} --- \texttt{uvx --from runops runo init \&\& source .venv/bin/activate}
  \item \textbf{Claude Code 起動} --- \texttt{claude}
  \item \textbf{以降はすべて AI に指示} --- 「環境を確認して」「ケースを作って」「サーベイを設計して」\ldots{}
\end{enumerate}

\begin{infobox}
Step 1--6 が唯一ユーザーが直接コマンドを実行する部分です (初回のみ)。Step 7 以降は Claude Code に自然言語で指示するだけで、runops のコマンドを覚える必要はありません。
\end{infobox}

% ====== Section 10 ======
\section*{\color{darkblue} 10.\quad 参考リンク}

\textbf{runops:}
\begin{itemize}[nosep]
  \item GitHub: \url{https://github.com/Nkzono99/runops}
  \item ドキュメント: \texttt{tools/runops/docs/} (プロジェクト内)
\end{itemize}

\textbf{Claude Code:}
\begin{itemize}[nosep]
  \item 公式: \url{https://claude.ai/code}
  \item ドキュメント: \url{https://docs.anthropic.com/en/docs/claude-code}
  \item GitHub: \url{https://github.com/anthropics/claude-code}
\end{itemize}

\textbf{camphor3:}
\begin{itemize}[nosep]
  \item マニュアル: \url{https://web.kudpc.kyoto-u.ac.jp/manual/ja}
\end{itemize}

\end{document}
